\documentclass[aspectratio=169]{beamer}
% Shared Beamer configuration for MUS2640 lecture decks (included after \documentclass).
\usepackage[utf8]{inputenc}
\usepackage[T1]{fontenc}
\usepackage{lmodern}
\usepackage{graphicx}
\usepackage{booktabs}

\usetheme{Madrid}
\usecolortheme{dolphin}
\usefonttheme{professionalfonts}

\setbeamercovered{transparent}
\setbeamertemplate{navigation symbols}{}
\setbeamertemplate{footline}[frame number]

\newcommand{\coursename}{MUS2640 · Sensing Sound and Music}
\newcommand{\instituteuo}{University of Oslo · Department of Musicology / RITMO}


\title{Week 8 · The body}
\subtitle{\coursename}
\institute{\instituteuo}
\date{}

\begin{document}

\begin{frame}
  \titlepage
\end{frame}

\begin{frame}{Aims}
  \begin{itemize}
    \item Embodied music cognition: perception--action loops; gesture.
    \item \textbf{4E} sketch: embodied, embedded, enactive, extended.
    \item Joint action, dance, and tactile/haptic channels (overview).
    \item Motion capture paradigms — qualitative vs quantitative.
  \end{itemize}
\end{frame}

\begin{frame}{Roadmap}
  \begin{enumerate}
    \item Performer vs perceiver motion; sound-producing vs sound-accompanying gestures.
    \item Involuntary motion: micromotion and the standstill studies (mm/s, not visible on video).
    \item Anatomy/biomechanics — enough to read motion studies (high level).
    \item Joint ensembles / synchrony; dance--music relations (cultural specificity).
    \item Methods: video, MoCap, sensors — trade-offs.
  \end{enumerate}
\end{frame}

\begin{frame}{Multimodality reminder}
  \begin{itemize}
    \item Seeing performance reshapes heard timing and expression (Vision week hooks).
    \item Vibration and touch carry pulse when audition is masked or inaccessible.
  \end{itemize}
\end{frame}

\begin{frame}{In-class}
  \begin{itemize}
    \item Observe or film a short performance clip — list visible gestures tied to phrasing.
    \item Optional: simple movement observation protocol (pairs).
  \end{itemize}
\end{frame}

\begin{frame}{Outside class}
  \begin{itemize}
    \item Reading + short plan for an essay or reflection (per course schedule).
  \end{itemize}
\end{frame}

\begin{frame}{Summary}
  \begin{itemize}
    \item Music cognition ``in the wild'' is embodied and socially coupled.
    \item Next: \textbf{Physiology} — autonomic measures (HR, EDA, breathing) and ethics of measurement.
  \end{itemize}
\end{frame}

\end{document}
